% !PRINTED 217
\weberpage{217}

An invariant of degree $n(n-1)$ is also the discriminant $\Delta$ of equation (3), or the square of the difference product

\weq{4}{\begin{gathered} \Pi = (y_1 - y_2)(y_1 - y_3)\cdots \\ (y_2-y_3)\cdots \end{gathered}}

To obtain closer insight into the formation of this quantity, we recall that we obtain $y$ when, in

\[ \frac{f(t)}{t-x} - \frac{1}{n}f'(t) \]

we replace $t^k$ by $t_k$. We therefore obtain $y_1-y_2$ from the formula

\weq{5}{y_1-y_2 = (x_1-x_2)\frac{f(t)}{(t-x_1)(t-x_2)},}

when we make the same replacement.

The quotient

\[ \frac{f(t)}{(t-x_i)(t-x_k)} \]

is an integral rational function of $t$ of degree $n-2$. If we replace $t^k$ in it by $t_k$, let it pass over into $t_{i,k}$; we then have

\weq{6}{y_i - y_k = (x_i-x_k)t_{i,k}.}

If we therefore set

\weq{7}{\begin{gathered} a_0^{n-1}\Theta = t_{1,2}t_{1,3}\ldots t_{1,n} \\ t_{2,3}\ldots t_{2,n} \\ t_{n-1,n} \end{gathered}}

then $\Theta$ is a homogeneous integral function of degree $\frac{1}{2}n(n-1)$ with respect to the $t$, and moreover, as a symmetric function of the $x$, it can be expressed rationally in terms of the coefficients $a$.

But from (6) it follows, if $D$, as in \S.\,46, (3), denotes the discriminant of $f(x)$, that

\weq{8}{\Delta = D\Theta^2,}

from which it is to be concluded that $\Theta$ is an invariant \gesperrt{which no longer contains the denominator $a_0$}. $\Delta$ is of degree $n(n-1)$ in the coefficients $a$, while $D$ by \S.\,46 is of degree $2n-2$; from this it follows that $\Theta$ is of degree $\frac{1}{2}(n-1)(n-2)$ in the coefficients $a$.
